% Intended LaTeX compiler: pdflatex
\documentclass[a4paper,12pt]{article}
\usepackage[utf8]{inputenc}
\usepackage[T1]{fontenc}
\usepackage{graphicx}
\usepackage{longtable}
\usepackage{wrapfig}
\usepackage{rotating}
\usepackage[normalem]{ulem}
\usepackage{amsmath}
\usepackage{amssymb}
\usepackage{capt-of}
\usepackage{hyperref}
\usepackage{\string~/.emacs.d/common}
\author{mahmood}
\date{\today}
\title{deduction theorem}
\hypersetup{
 pdfauthor={mahmood},
 pdftitle={deduction theorem},
 pdfkeywords={},
 pdfsubject={},
 pdfcreator={Emacs 30.0.50 (Org mode 9.6.6)}, 
 pdflang={English}}
\begin{document}

\maketitle
\begin{note}
this document and others can be found on my website\\
\url{https://mahmoodsheikh36.github.io/}
\end{note}
in a \textbf{complete} \href{20230524203100-formal_system.org}{formal system}, in order to prove \(\alpha \to \beta\) we can add \(\alpha\) to the set of assumptions and then prove \(\beta\)\\[0pt]
a special case of the \textbf{deduction theorem} is: \(p \vdash q\) iff \(\vdash (p \to q)\)\\[0pt]
when writing a \href{20230525071323-formal_proof.org}{formal proof}, we write \(IP\) (atleast in my college) to denote the usage of deduction theorem\\[0pt]
\begin{my_example}
infer \((\lnot P \to S) \to S\) from \(P \to S\)\\[0pt]
\begin{align}
  P \to S &\quad P_1\\
  \lnot P \to S &\quad IP\\
  \lnot S \to P &\quad 2,E21\\
  \lnot S \to S &\quad 1,3,I13\\
  \lnot (\lnot S) \lor S &\quad 4,E20\\
  S \lor S &\quad 5,E9\\
  S &\quad 5,E5
\end{align}
\end{my_example}
\end{document}